\documentclass[11pt,a4paper]{article}
\usepackage[utf8]{inputenc}
\usepackage[T1]{fontenc}
\usepackage{lmodern}
\usepackage[french]{babel}
\usepackage{amsmath,amssymb,mathtools}
\usepackage{icomma}
\usepackage[version=4]{mhchem}
\usepackage{siunitx}
\sisetup{output-decimal-marker={,},per-mode=symbol}
\usepackage{xcolor}
\usepackage{colortbl}
\usepackage{tikz}
\usepackage[most]{tcolorbox}
\usepackage{enumitem}
\usepackage{array}
\usepackage{booktabs}
\usepackage{fancyhdr}
\usepackage[hidelinks]{hyperref}
\usetikzlibrary{arrows.meta,positioning,calc}
\usepackage[a4paper,margin=2.1cm,headheight=26pt,headsep=10pt]{geometry}

\definecolor{primary}{RGB}{146,95,161}
\definecolor{primarydark}{RGB}{92,55,108}
\definecolor{acidecol}{RGB}{205,60,45}
\definecolor{basecol}{RGB}{40,110,200}
\definecolor{excol}{RGB}{0,135,90}
\definecolor{methcol}{RGB}{90,90,100}

\tcbset{boxbase/.style={enhanced,breakable,boxrule=0.9pt,arc=2.5pt,left=3mm,right=3mm,top=2.3mm,bottom=2.3mm,fonttitle=\bfseries,coltitle=white,toptitle=0.6mm,bottomtitle=0.6mm}}
\newtcolorbox[auto counter]{exobox}[1][]{boxbase,colback=primary!4,colframe=primary,title={\textsc{Exercice}~\thetcbcounter\ifx\relax#1\relax\relax\else\textmd{ --- \itshape #1}\fi}}
\newtcolorbox{consigne}{boxbase,colback=methcol!7,colframe=methcol,sharp corners}

\pagestyle{fancy}
\fancyhf{}
\fancyhead[L]{\small\textcolor{primarydark}{\textbf{Fiche 10} — Thème 1 : couples acide/base}}
\fancyhead[R]{\small\textcolor{primarydark}{\textsc{Exercices — Facile}}}
\fancyfoot[C]{\small\thepage}
\renewcommand{\headrulewidth}{0.8pt}
\renewcommand{\headrule}{\hbox to\headwidth{\color{primary}\leaders\hrule height \headrulewidth\hfill}}

\setlength{\parindent}{0pt}
\setlength{\parskip}{3pt}

\begin{document}

\begin{center}
{\LARGE\bfseries\color{primarydark} Thème 1 --- Niveau Facile}\\[2pt]
{\color{primary}\itshape Couples acide/base conjugués : le geste $\pm\,\ce{H+}$}
\end{center}

\begin{consigne}
\textbf{Objectif du niveau.} Automatiser la construction d'un couple : passer d'un acide à sa base conjuguée
(on enlève un \ce{H+}) et d'une base à son acide conjugué (on ajoute un \ce{H+}), sans jamais se tromper de
charge. Aucun calcul : uniquement des formules chimiques.
\end{consigne}

\begin{exobox}[donner la base conjuguée d'un acide]
Pour chacun des acides suivants, écrire la \textbf{base conjuguée} (on retire un \ce{H+}), en soignant la charge :
\begin{enumerate}[label=\alph*),leftmargin=1.6em,itemsep=1pt]
  \item \ce{HCl}
  \item \ce{CH3COOH} (acide acétique)
  \item \ce{NH4+} (ion ammonium)
  \item \ce{H2CO3} (acide carbonique)
  \item \ce{HSO4-} (ion hydrogénosulfate)
\end{enumerate}
\end{exobox}

\begin{exobox}[donner l'acide conjugué d'une base]
Pour chacune des bases suivantes, écrire l'\textbf{acide conjugué} (on ajoute un \ce{H+}), en soignant la charge :
\begin{enumerate}[label=\alph*),leftmargin=1.6em,itemsep=1pt]
  \item \ce{OH-} (ion hydroxyde)
  \item \ce{NH3} (ammoniac)
  \item \ce{CO3^{2-}} (ion carbonate)
  \item \ce{HCO3-} (ion hydrogénocarbonate)
  \item \ce{F-} (ion fluorure)
\end{enumerate}
\end{exobox}

\begin{exobox}[qui donne, qui capte ?]
Dans chacune des réactions ci-dessous, désigner \textbf{l'espèce qui joue le rôle d'acide} (celle qui cède
le \ce{H+}) et \textbf{celle qui joue le rôle de base} (celle qui capte le \ce{H+}).
\begin{enumerate}[label=\alph*),leftmargin=1.6em,itemsep=2pt]
  \item $\ce{HCl + H2O -> Cl- + H3O+}$
  \item $\ce{NH3 + H2O <=> NH4+ + OH-}$
  \item $\ce{CH3COOH + H2O <=> CH3COO- + H3O+}$
\end{enumerate}
\end{exobox}

\begin{exobox}[écrire la demi-équation d'un couple]
Écrire la demi-équation acide-base $\ce{acide <=> base + H+}$ associée à chacun des couples suivants :
\begin{enumerate}[label=\alph*),leftmargin=1.6em,itemsep=1pt]
  \item couple \ce{HNO2}/\ce{NO2-}
  \item couple \ce{NH4+}/\ce{NH3}
  \item couple \ce{HCO3-}/\ce{CO3^{2-}}
\end{enumerate}
\end{exobox}

\begin{exobox}[acide ou base de Lewis ?]
Selon Lewis, un \textbf{acide} accepte un doublet d'électrons (atome à orbitale vide) et une \textbf{base}
donne un doublet d'électrons. Classer chaque espèce en « acide de Lewis » ou « base de Lewis » :
\begin{enumerate}[label=\alph*),leftmargin=1.6em,itemsep=1pt]
  \item \ce{NH3} (l'azote porte un doublet non liant)
  \item \ce{BF3} (le bore possède une orbitale vide)
  \item \ce{H+}
  \item \ce{H2O} (l'oxygène porte deux doublets non liants)
\end{enumerate}
\end{exobox}

\end{document}
